\subsection{Divisor classes attached to finitely generated modules}

We keep the notation and hypotheses of nos. 2 to 6. Recall that C(A) (or simply C) denotes the divisor class group of A, the quotient of D(A) by the subgroup of principal divisors. For every divisor \(d \in D\) , we shall denote by \(c(d)\) its class in C.

PROPOSITION 15. Let M be a finitely generated A-module. There exists a free submodule L of M such that M/L is a torsion module and the element \(c(\chi(M/L))\) of C does not depend on the free submodule L chosen.

We write \(S = A - \{0\}\) and let \(V = S^{-1}M = M \otimes_{A} K\) ; if \(n\) is the rank of \(V\) over \(K\) , there exist \(n\) elements \(e_i (1 \leqslant i \leqslant n)\) of \(M\) whose canonical images in \(V\) form a basis of \(V\) ; these elements are obviously linearly independent in \(M\) and hence generate a free submodule \(L\) of \(M\) such that \(S^{-1}(M/L) = S^{-1}M/S^{-1}L = 0\) , so that \(M/L\) is a torsion module.

Now let \(L_{1}\) be another free submodule of M of rank n.~Since \(S^{-1}L = S^{-1}L_{1}\) , there exists \(s \in S\) such that \(sL_{1} \subset L\) ; we may therefore limit ourselves to proving that, if \(L_{1} \subset L_{2}\) are two free submodules of M of rank n, then

\[
c (\chi (\mathbf {M} / \mathbf {L} _ {1})) = c (\chi (\mathbf {M} / \mathbf {L} _ {2})).
\]

Now, \(\chi (\mathbf{M} / \mathbf{L}_1) = \chi (\mathbf{M} / \mathbf{L}_2) + \chi (\mathbf{L}_2 / \mathbf{L}_1)\) and it follows from no. 6, Corollary to Proposition 14 that \(\chi (\mathbf{L}_2 / \mathbf{L}_1)\) is a principal divisor and therefore

\[
c (\chi (\mathbf {L} _ {2} / \mathbf {L} _ {1})) = 0.
\]

The element \(c(\chi(\mathbf{M} / \mathbf{L}))\) will be denoted by \(-c(\mathbf{M})\) in what follows; we shall say that \(c(\mathbf{M})\) is the divisor class attached to \(\mathbf{M}\) .

PROPOSITION 16. (i) Let \(0 \to \mathbf{M}_1 \xrightarrow{f} \mathbf{M}_2 \xrightarrow{g} \mathbf{M}_3 \to 0\) be an exact sequence of finitely generated A-modules. Then

\[
c (\mathbf {M _ {2}}) = c (\mathbf {M _ {1}}) + c (\mathbf {M _ {3}}).
\]

\begin{enumerate}
\def\labelenumi{(\roman{enumi})}
\setcounter{enumi}{1}
\item
  If there exists a pseudo-isomorphism from \(\mathbf{M}_1\) to \(\mathbf{M}_2\) , then \(c(\mathbf{M}_1) = c(\mathbf{M}_2)\) .
\item
  If \(\mathbf{T}\) is a torsion module, then \(c(\mathbf{T}) = -c(\chi(\mathbf{T}))\) .
\item
  If \(\mathfrak{a} \neq 0\) is a fractional ideal of \(\mathbf{K}\) , then
\end{enumerate}

\[
c (\mathfrak {a}) = c (\operatorname{div} (\mathfrak {a})).
\]

\begin{enumerate}
\def\labelenumi{(\alph{enumi})}
\setcounter{enumi}{21}
\tightlist
\item
  If \(\mathbf{L}\) is a free A-module, then \(c(\mathbf{L}) = 0\) .
\end{enumerate}

To prove (i), consider a free sub-module \(\mathbf{L}_1\) (resp. \(\mathbf{L}_3\) ) of \(\mathbf{M}_1\) (resp. \(\mathbf{M}_3\) ) such that \(\mathbf{M}_1 / \mathbf{L}_1\) (resp. \(\mathbf{M}_3 / \mathbf{L}_3\) ) is a torsion module. Since \(\mathbf{L}_3\) is free and \(g\) is surjective, there exists in \(\overline{g}^{-1} (\mathbf{L}_3)\) a free complement \(\mathbf{L}_{23}\) of \(\operatorname {Ker}(g)\) which is isomorphic to \(\mathbf{L}_3\) (Algebra, Chapter II, § 1, no. 11, Proposition 21); but \(\operatorname {Ker}(g) = f(\mathbf{M}_1)\)

contains \(f(\mathbf{L}_1) = \mathbf{L}_{12}\) which is free since \(f\) is injective. The sum \(\mathbf{L}_2 = \mathbf{L}_{12} + \mathbf{L}_{23}\) is direct and \(\mathbf{L}_2\) is therefore a free submodule of \(\mathbf{M}_2\) . There is moreover the commutative diagram:

\[
\begin{array}{c} 0 \longrightarrow \mathrm{L} _ {1} \longrightarrow \mathrm{L} _ {2} \longrightarrow \mathrm{L} _ {3} \longrightarrow 0 \\ \Biggl \downarrow \qquad \qquad \qquad \Biggl \downarrow \qquad \qquad \Biggl \downarrow \\ 0 \longrightarrow \mathrm{M} _ {1} \xrightarrow [ f ]{} \mathrm{M} _ {2} \xrightarrow [ g ]{} \mathrm{M} _ {3} \longrightarrow 0 \end{array}
\]

where the rows are exact and the vertical arrows are injections. We therefore obtain from the snake diagram (Chapter I, § 1, no. 4, Proposition 2) the exact sequence:

\[
0 \rightarrow \mathrm{M} _ {1} / \mathrm{L} _ {1} \rightarrow \mathrm{M} _ {2} / \mathrm{L} _ {2} \rightarrow \mathrm{M} _ {3} / \mathrm{L} _ {3} \rightarrow 0.
\]

As \(\mathbf{M}_1 / \mathbf{L}_1\) and \(\mathbf{M}_3 / \mathbf{L}_3\) are torsion modules, this exact sequence shows first that so is \(\mathbf{M}_2 / \mathbf{L}_2\) and then, by virtue of Proposition 10 of no. 5, that:

\[
\chi (\mathrm{M} _ {2} / \mathrm{L} _ {2}) = \chi (\mathrm{M} _ {1} / \mathrm{L} _ {1}) + \chi (\mathrm{M} _ {3} / \mathrm{L} _ {3})
\]

which proves (i).

Assertions (iii) and (v) are obvious from the definition. We prove (ii). Therefore let \(f \colon \mathbf{M}_1 \to \mathbf{M}_2\) be a pseudo-isomorphism and let \(\mathbf{L}_1\) be a free submodule of \(\mathbf{M}_1\) such that \(\mathbf{M}_1 / \mathbf{L}_1\) is a torsion module. We set \(\mathbf{L}_2 = f(\mathbf{L}_1)\) ; as \(\operatorname{Ker}(f)\) is pseudo-zero, it is a torsion module, hence \(\operatorname{Ker}(f) \cap \mathbf{L}_1 = 0\) and therefore \(\mathbf{L}_2\) is free. Let \(\bar{f} \colon \mathbf{M}_1 / \mathbf{L}_1 \to \mathbf{M}_2 / \mathbf{L}_2\) be the homomorphism derived from \(f\) by taking quotients; \(\operatorname{Ker}(\bar{f})\) is isomorphic to \(\operatorname{Ker}(f)\) and \(\operatorname{Coker}(\bar{f})\) to \(\operatorname{Coker}(f)\) and hence \(\bar{f}\) is a pseudo-isomorphism; moreover

\[
\operatorname{Coker} (\bar {f}) = \mathrm{M} _ {2} / f (\mathrm{M} _ {1})
\]

is a torsion module and so is \(f(\mathbf{M}_1) / \mathbf{L}_2 = \bar{f} (\mathbf{M}_1 / \mathbf{L}_1)\) , hence \(\mathbf{M}_2 / \mathbf{L}_2\) is a torsion module and it follows from no. 5, Proposition 10 (ii) that

\[
\chi (\mathrm{M} _ {1} / \mathrm{L} _ {1}) = \chi (\mathrm{M} _ {2} / \mathrm{L} _ {2}).
\]

Finally it remains to prove (iv). Let \(x \in \mathbf{K}^*\) be such that \(a \subset xA\) . By considering the exact sequence \(0 \to a \to xA \to xA / a \to 0\) , we obtain

\[
c (\mathfrak {a}) = c (x \mathrm{A}) - c (x \mathrm{A} / \mathfrak {a}) = - c (x \mathrm{A} / \mathfrak {a})
\]

by (i) and (v). But \(x\mathrm{A} / \mathfrak{a}\) is isomorphic to \(\mathrm{A} / x^{-1}\mathfrak{a}\) , whence, by virtue of (iii),

\[
c (x \mathrm{A} / \mathrm{a}) = - c (\chi (\mathrm{A} / x ^ {- 1} \mathrm{a})) = - c (\operatorname{div} (x ^ {- 1} \mathrm{a})) = - c (\operatorname{div} (\mathrm{a}))
\]

(no. 5, Proposition 12). This completes the proof.

When M is a lattice of V with respect to A, \(\chi(\mathrm{M}/\mathrm{L}) = -\chi(\mathrm{M}, \mathrm{L})\) (no. 6, Proposition 14); let \((e_{i})_{1 \leqslant i \leqslant n}\) be a basis of L, \(e = e_{1} \wedge e_{2} \wedge \cdots \wedge e_{n}\) and

\(M_{w} = a.e\) (notation of no. 6); then \(\chi(M, L) = \text{div}(a)\) , whence \(c(M) = c(\text{div}(a))\) , which generalizes Proposition 16 (v).

COROLLARY 1. Let \(0 \to \mathbf{M}_n \xrightarrow{u} \mathbf{M}_{n-1} \to \cdots \to \mathbf{M}_0 \to 0\) be an exact sequence of finitely generated A-modules. Then

\[
\sum_ {i = 0} ^ {n} (- 1) ^ {i} c (\mathbf {M} _ {i}) = 0.
\]

We argue by induction on \(n\) , the case \(n = 2\) being Proposition 16 (i). If \(\mathbf{M}_{n-1}' = \operatorname{Coker}(u)\) , there are the two exact sequences:

\[
0 \rightarrow \mathrm{M} _ {n} \rightarrow \mathrm{M} _ {n - 1} \rightarrow \mathrm{M} _ {n - 1} ^ {\prime} \rightarrow 0
\]

\[
0 \rightarrow \mathrm{M} _ {n - 1} ^ {\prime} \rightarrow \mathrm{M} _ {n - 2} \rightarrow \dots \rightarrow \mathrm{M} _ {0} \rightarrow 0.
\]

The first shows that \(\mathbf{M}_{n - 1}^{\prime}\) is finitely generated and the induction hypothesis gives

\[
(- 1) ^ {n - 1} c (\mathbf {M} _ {n - 1} ^ {\prime}) + \sum_ {i = 0} ^ {n - 2} (- 1) ^ {i} c (\mathbf {M} _ {i}) = 0
\]

and

\[
c (\mathbf {M} _ {n - 1} ^ {\prime}) = c (\mathbf {M} _ {n - 1}) - c (\mathbf {M} _ {n}),
\]

whence the corollary.

An exact sequence

\[
0 \rightarrow \mathrm{L} _ {n} \rightarrow \mathrm{L} _ {n - 1} \rightarrow \dots \rightarrow \mathrm{L} _ {0} \rightarrow \mathrm{E} \rightarrow 0
\]

where the \(L_{i}\) ( \(0 \leqslant i \leqslant n\) ) are finitely generated free A-modules, is called a finite free resolution of the A-module E.

COROLLARY 2. If a divisorial fractional ideal \(\mathfrak{a} \neq 0\) of A admits a finite free resolution, it is principal.

In fact we apply Corollary 1 to a finite free resolution of \(\mathfrak{a}\) :

\[
0 \rightarrow \mathrm{L} _ {n} \rightarrow \mathrm{L} _ {n - 1} \rightarrow \dots \rightarrow \mathrm{L} _ {0} \rightarrow \mathfrak {a} \rightarrow 0.
\]

By virtue of Proposition 16 (v), \(c(a) = 0\) and hence, by virtue of Proposition 16 (iv), \(\text{div}(a)\) is principal; as \(a\) is assumed to be divisorial, it is principal (§ 1, no. 1).

COROLLARY 3. If every divisorial ideal \(\neq 0\) of A admits a finite free resolution, A is factorial.

This is an immediate consequence of Corollary 2 and § 3, no. 1, Definition 1.

* We shall see later that a regular local ring satisfies the hypothesis of Corollary

3 and therefore is factorial. *

If M is a finitely generated A-module, we shall denote its rank by \(r(\mathbf{M})\) (recall that it is the rank over K of \(\mathbf{M}_{(\mathbf{K})} = \mathbf{M} \otimes_{\mathbf{A}} \mathbf{K}\) ); if \(0 \to M_{1} \to M_{2} \to M_{3} \to 0\) is an exact sequence of finitely generated A-modules, the sequence

\[
0 \rightarrow (\mathrm{M} _ {1}) _ {(\mathbf {K})} \rightarrow (\mathrm{M} _ {2}) _ {(\mathbf {K})} \rightarrow (\mathrm{M} _ {3}) _ {(\mathbf {K})} \rightarrow 0
\]

is also exact and hence \(r(\mathbf{M}_2) = r(\mathbf{M}_1) + r(\mathbf{M}_3)\) . We write

\[
\gamma (\mathbf {M}) = (r (\mathbf {M}), c (\mathbf {M})) \in \mathbf {Z} \times \mathbf {C} (\mathbf {A});
\]

γ therefore satisfies property (i) of Proposition 16 and, if M is pseudo-zero, γ(M) = 0 (since M is a torsion module). There exists a unique mapping from F(A) to Z × C(A), also denoted by γ, such that γ(M) = γ(cl(M)) for every finitely generated A-module M. We shall see that the above properties essentially characterize γ:

PROPOSITION 17. Let G be a commutative group written additively and \(\phi\) a mapping from the set \(F(A)\) of classes of finitely generated A-modules to G; for every finitely generated A-module M we also write, by an abuse of language, \(\phi(M) = \phi(cl(M))\) . Suppose the following conditions are satisfied:

\begin{enumerate}
\def\labelenumi{(\arabic{enumi})}
\item
  If \(0 \to \mathbf{M}_1 \to \mathbf{M}_2 \to \mathbf{M}_3 \to 0\) is an exact sequence of finitely generated A-modules, then \(\phi(\mathbf{M}_2) = \phi(\mathbf{M}_1) + \phi(\mathbf{M}_3)\) .
\item
  If \(\mathbf{T}\) is pseudo-zero, then \(\phi (\mathbf{T}) = 0\) .
\end{enumerate}

Then there exists a unique homomorphism \(\theta: \mathbf{Z} \times \mathbf{C} \to \mathbf{G}\) such that \(\phi = \theta \circ \gamma\) .

By virtue of Proposition 16 (iv), every element of \(\mathbf{Z} \times \mathbf{C}\) is of the form \((r(\mathbf{M}), c(\mathbf{M}))\) for some suitable finitely generated A-module M; whence the uniqueness of \(\theta\) . We apply Proposition 11 of no. 5 to the restriction of \(-\phi\) to \(T(\mathbf{A})\) : then there exists a homomorphism \(\theta_0: \mathbf{D} \to \mathbf{G}\) such that

\[
- \phi (\mathbf {T}) = \theta_ {0} (\chi (\mathbf {T}))
\]

for every finitely generated torsion A-module T. Let \(x\) be a non-zero element of A; applying property (1) to the exact sequence:

\[
0 \longrightarrow \mathrm{A} \xrightarrow {h _ {x}} \mathrm{A} \longrightarrow \mathrm{A} / x \mathrm{A} \longrightarrow 0
\]

where \(h_{x}\) is multiplication by x, we obtain \(\phi(\mathbf{A}/x\mathbf{A}) = 0\) , whence \(\theta_{0}(\mathrm{div}(x)) = 0\) . Taking quotients, \(\theta_{0}\) therefore defines a homomorphism \(\theta_{1}: C \to G\) and \(\phi(T) = \theta_{1}(c(T))\) for every torsion A-module T. We show now that the homomorphism \(\theta\) defined by \(\theta(n, z) = n.\phi(\mathbf{A}) + \theta_{1}(z)\) solves the problem. For this, we write \(\phi'(\mathbf{M}) = \phi(\mathbf{M}) - \theta(\gamma(\mathbf{M}))\) for every finitely generated A-module M; clearly condition (1) is still satisfied if \(\phi\) is replaced by \(\phi'\) . Moreover, \(\phi'(\mathbf{M}) = 0\) when M is a torsion module or a free module; but as for every finitely generated A-module M, there exists a free sub-module L of M such that M/L is a torsion module (Proposition 15), property (1) shows that \(\phi'(\mathbf{M}) = 0\) for every finitely generated A-module M.

* In the language of Abelian categories, Proposition 17 shows that \(\mathbf{Z} \times \mathrm{C}(\mathbf{A})\) is canonically isomorphic to the Grothendieck group of the quotient category \(\mathcal{F}/\mathcal{F}'\) , where \(\mathcal{F}\) is the category of finitely generated A-modules and \(\mathcal{F}'\) the full sub-category of \(\mathcal{F}\) consisting of the pseudo-zero modules. *

